
\documentclass{article}
\usepackage{colortbl}
\usepackage{makecell}
\usepackage{multirow}
\usepackage{supertabular}

\begin{document}

\newcounter{utterance}

\centering \large Interaction Transcript for game `clean\_up', experiment `2\_hard\_5obj\_sl', episode 3 with deepseek{-}v4{-}pro.
\vspace{24pt}

{ \footnotesize  \setcounter{utterance}{1}
\setlength{\tabcolsep}{0pt}
\begin{supertabular}{c@{$\;$}|p{.15\linewidth}@{}p{.15\linewidth}p{.15\linewidth}p{.15\linewidth}p{.15\linewidth}p{.15\linewidth}}
    \# & $\;$A & \multicolumn{4}{c}{Game Master} & $\;\:$B\\
    \hline

    \theutterance \stepcounter{utterance}  
    & & \multicolumn{4}{p{0.6\linewidth}}{
        \cellcolor[rgb]{0.9,0.9,0.9}{
            \makecell[{{p{\linewidth}}}]{
                \texttt{\tiny{[P1$\langle$GM]}}
                \texttt{Jaz sem vaš vodja igre, vi pa igrate skupinsko igro z naslednjo mrežo kot igralno ploščo:} \\
\\ 
\texttt{```} \\
\texttt{ 1234567} \\
\texttt{╔═╤═╤═══╗} \\
\texttt{║C│◌│◌TI║ 1} \\
\texttt{║◌├─┤W◌◌║ 2} \\
\texttt{╟─┤◌├─┬─╢ 3} \\
\texttt{║◌│◌│◌│◌║ 4} \\
\texttt{╟─┼─┘◌│◌║ 5} \\
\texttt{║◌│◌H┌┴─╢ 6} \\
\texttt{║◌│◌◌│◌◌║ 7} \\
\texttt{╚═╧══╧══╝} \\
\\ 
\texttt{```} \\
\\ 
\texttt{* Na zgornjem robu so prikazane koordinate x, ki naraščajo proti desni, na desnem robu pa koordinate y, ki naraščajo navzdol.} \\
\texttt{* Na mrežo so naključno postavljeni naslednji predmeti: 'W', 'I', 'T', 'C', 'H'.} \\
\\ 
\\ 
\texttt{Drugi igralec vidi različico igralne plošče, na kateri so predmeti postavljeni na različna naključna mesta. Vi ne morete videti plošče drugega igralca, on pa ne more videti vaše.} \\
\\ 
\\ 
\texttt{**Cilj:**} \\
\\ 
\texttt{Oba igralca morata premikati predmete na svojem ozadju tako, da se enaki predmeti znajdejo na istih koordinatah. Z drugim igralcem se morate dogovoriti o skupni ciljni konfiguraciji.} \\
\\ 
\\ 
\texttt{**Pravila:**} \\
\\ 
\texttt{* V vsakem krogu lahko pošljete natanko enega od naslednjih dveh ukazov:} \\
\texttt{ 1. `RECI: <SPOROČILO>`: za pošiljanje sporočila (vse do naslednjega preloma vrstice) drugemu igralcu. Posredoval ga bom vašemu partnerju.} \\
\texttt{2. `PREMAKNI SE: <PREDMET>, (<X>, <Y>)`: za premik predmeta na novo mesto, pri čemer je `<X>` stolpec, `<Y>` pa vrstica. Obvestil vas bom, ali je bila vaša premestitev uspešna ali ne.} \\
\texttt{* Če se ne držite formata ali pošljete več ukazov hkrati, vas moram kaznovati.} \\
\texttt{* Če oba igralca zbereta več kot 12 kazni, oba izgubita igro.} \\
\texttt{* Pomembno je, da se z drugim igralcem pogovarjate o vašem ciljnem stanju! Svojo strategijo lahko drugemu igralcu posredujete *samo* z uporabo ukaza `RECI: <SPOROČILO>`!} \\
\\ 
\\ 
\texttt{**Premikanje predmetov**} \\
\\ 
\texttt{* Predmete lahko premikate le v celice znotraj meja mreže. Ciljna celica mora biti prazna, tj. vsebovati mora samo simbol '◌'.} \\
\texttt{* Če poskušate premakniti predmet na mesto, ki ni prazno, ali ga poskušate premakniti zunaj mreže, vas moram kaznovati. Dobili boste še en poskus.} \\
\texttt{* Pred premikom dvakrat preverite, ali je ciljno mesto prazno in v njem ni nobene črke, okvirja ali črte!} \\
\\ 
\\ 
\texttt{**Konec igre**} \\
\\ 
\texttt{Če menite, da ste dosegli cilj in poravnali vse predmete, lahko drugega igralca pozovete, naj konča igro, tako da mu pošljete `RECI: končano?`. Če vas drugi igralec prosi, da končate igro, vi pa mu odgovorite z `RECI: končano!`, se igra konča.} \\
\\ 
\texttt{Oba igralca zmagata, če se igra konča v 20 krogih, pri čemer je en krog definiran tako, da dva igralca pošljeta vsak po en veljaven ukaz.} \\
\\ 
\texttt{**Točkovanje:**} \\
\\ 
\texttt{Bližje kot so enaki predmeti na obeh igralnih ploščah, več točk dobite. Kazni zmanjšajo število točk. Ali lahko premagate rekord?} \\
\\ 
\\ 
\texttt{Pošljite sporočilo drugemu igralcu, da začnete igro!} \\
            }
        }
    }
    & \\ \\

    \theutterance \stepcounter{utterance}  
    & \multicolumn{4}{p{0.6\linewidth}}{
        \cellcolor[rgb]{0.9,0.9,0.9}{
            \makecell[{{p{\linewidth}}}]{
                \texttt{\tiny{[P1$\rangle$GM]}}
                \texttt{RECI: Pozdravljen! Začniva igro. Na moji plošči so predmeti na naslednjih položajih: C je na (1,1), T je na (5,1), I je na (6,1), W je na (4,2), H je na (4,6). Kje so predmeti na tvoji plošči?} \\
            }
        }
    }
    & & \\ \\

    \theutterance \stepcounter{utterance}  
    & & & \multicolumn{2}{p{0.3\linewidth}}{
        \cellcolor[rgb]{0.9,0.9,0.9}{
            \makecell[{{p{\linewidth}}}]{
                \texttt{\tiny{[GM$|$GM]}}
                \texttt{RECI: Pozdravljen! Začniva igro. Na moji plošči so predmeti na naslednjih položajih: C je na (1,1), T je na (5,1), I je na (6,1), W je na (4,2), H je na (4,6). Kje so predmeti na tvoji plošči?} \\
            }
        }
    }
    & & \\ \\

    \theutterance \stepcounter{utterance}  
    & & \multicolumn{4}{p{0.6\linewidth}}{
        \cellcolor[rgb]{0.9,0.9,0.9}{
            \makecell[{{p{\linewidth}}}]{
                \texttt{\tiny{[GM$\rangle$P2]}}
                \texttt{Jaz sem vaš vodja igre, vi pa igrate skupinsko igro z naslednjo mrežo kot igralno ploščo:} \\
\\ 
\texttt{```} \\
\texttt{ 1234567} \\
\texttt{╔═╤═╤═══╗} \\
\texttt{║H│◌│◌◌◌║ 1} \\
\texttt{║◌├─┤◌◌◌║ 2} \\
\texttt{╟─┤◌├─┬─╢ 3} \\
\texttt{║◌│◌│C│◌║ 4} \\
\texttt{╟─┼─┘◌│◌║ 5} \\
\texttt{║◌│◌◌┌┴─╢ 6} \\
\texttt{║T│WI│◌◌║ 7} \\
\texttt{╚═╧══╧══╝} \\
\\ 
\texttt{```} \\
\\ 
\texttt{* Na zgornjem robu so prikazane koordinate x, ki naraščajo proti desni, na desnem robu pa koordinate y, ki naraščajo navzdol.} \\
\texttt{* Na mrežo so naključno postavljeni naslednji predmeti: 'W', 'I', 'T', 'C', 'H'.} \\
\\ 
\\ 
\texttt{Drugi igralec vidi različico igralne plošče, na kateri so predmeti postavljeni na različna naključna mesta. Vi ne morete videti plošče drugega igralca, on pa ne more videti vaše.} \\
\\ 
\\ 
\texttt{**Cilj:**} \\
\\ 
\texttt{Oba igralca morata premikati predmete na svojem ozadju tako, da se enaki predmeti znajdejo na istih koordinatah. Z drugim igralcem se morate dogovoriti o skupni ciljni konfiguraciji.} \\
\\ 
\\ 
\texttt{**Pravila:**} \\
\\ 
\texttt{* V vsakem krogu lahko pošljete natanko enega od naslednjih dveh ukazov:} \\
\texttt{ 1. `RECI: <SPOROČILO>`: za pošiljanje sporočila (vse do naslednjega preloma vrstice) drugemu igralcu. Posredoval ga bom vašemu partnerju.} \\
\texttt{2. `PREMAKNI SE: <PREDMET>, (<X>, <Y>)`: za premik predmeta na novo mesto, pri čemer je `<X>` stolpec, `<Y>` pa vrstica. Obvestil vas bom, ali je bila vaša premestitev uspešna ali ne.} \\
\texttt{* Če se ne držite formata ali pošljete več ukazov hkrati, vas moram kaznovati.} \\
\texttt{* Če oba igralca zbereta več kot 12 kazni, oba izgubita igro.} \\
\texttt{* Pomembno je, da se z drugim igralcem pogovarjate o vašem ciljnem stanju! Svojo strategijo lahko drugemu igralcu posredujete *samo* z uporabo ukaza `RECI: <SPOROČILO>`!} \\
\\ 
\\ 
\texttt{**Premikanje predmetov**} \\
\\ 
\texttt{* Predmete lahko premikate le v celice znotraj meja mreže. Ciljna celica mora biti prazna, tj. vsebovati mora samo simbol '◌'.} \\
\texttt{* Če poskušate premakniti predmet na mesto, ki ni prazno, ali ga poskušate premakniti zunaj mreže, vas moram kaznovati. Dobili boste še en poskus.} \\
\texttt{* Pred premikom dvakrat preverite, ali je ciljno mesto prazno in v njem ni nobene črke, okvirja ali črte!} \\
\\ 
\\ 
\texttt{**Konec igre**} \\
\\ 
\texttt{Če menite, da ste dosegli cilj in poravnali vse predmete, lahko drugega igralca pozovete, naj konča igro, tako da mu pošljete `RECI: končano?`. Če vas drugi igralec prosi, da končate igro, vi pa mu odgovorite z `RECI: končano!`, se igra konča.} \\
\\ 
\texttt{Oba igralca zmagata, če se igra konča v 20 krogih, pri čemer je en krog definiran tako, da dva igralca pošljeta vsak po en veljaven ukaz.} \\
\\ 
\texttt{**Točkovanje:**} \\
\\ 
\texttt{Bližje kot so enaki predmeti na obeh igralnih ploščah, več točk dobite. Kazni zmanjšajo število točk. Ali lahko premagate rekord?} \\
\\ 
\texttt{Drugi igralec je začel igro s tem sporočilom:} \\
\\ 
\texttt{"Pozdravljen! Začniva igro. Na moji plošči so predmeti na naslednjih položajih: C je na (1,1), T je na (5,1), I je na (6,1), W je na (4,2), H je na (4,6). Kje so predmeti na tvoji plošči?"} \\
\\ 
\texttt{Kakšen je vaš prvi ukaz?} \\
            }
        }
    }
    & \\ \\

    \theutterance \stepcounter{utterance}  
    & & & \multicolumn{4}{p{0.6\linewidth}}{
        \cellcolor[rgb]{0.9,0.9,0.9}{
            \makecell[{{p{\linewidth}}}]{
                \texttt{\tiny{[GM$\langle$P2]}}
                \texttt{RECI: Pozdravljen! Moji predmeti so: H na (1,1), T na (1,7), W na (3,7), I na (4,7), C na (4,4). Predlagam, da se dogovoriva za ciljne koordinate. Kaj praviš na to, da postaviva vse predmete v vrstico 7, kjer imam jaz že T, W, I? Konkretno: H(1,7), T(2,7), W(3,7), I(4,7), C(5,7)?} \\
            }
        }
    }
     \\ \\

    \theutterance \stepcounter{utterance}  
    & & & \multicolumn{2}{p{0.3\linewidth}}{
        \cellcolor[rgb]{0.9,0.9,0.9}{
            \makecell[{{p{\linewidth}}}]{
                \texttt{\tiny{[GM$|$GM]}}
                \texttt{RECI: Pozdravljen! Moji predmeti so: H na (1,1), T na (1,7), W na (3,7), I na (4,7), C na (4,4). Predlagam, da se dogovoriva za ciljne koordinate. Kaj praviš na to, da postaviva vse predmete v vrstico 7, kjer imam jaz že T, W, I? Konkretno: H(1,7), T(2,7), W(3,7), I(4,7), C(5,7)?} \\
            }
        }
    }
    & & \\ \\

    \theutterance \stepcounter{utterance}  
    & & \multicolumn{4}{p{0.6\linewidth}}{
        \cellcolor[rgb]{0.9,0.9,0.9}{
            \makecell[{{p{\linewidth}}}]{
                \texttt{\tiny{[P1$\langle$GM]}}
                \texttt{Vaše sporočilo je bilo posredovano drugemu igralcu.} \\
\\ 
\\ 
\texttt{Trenutno igrate krog 1 od največ 20.} \\
\texttt{Skupaj ste zbrali 0 od 12 kazni.} \\
\texttt{Drugi igralec je poslal to sporočilo:} \\
\\ 
\\ 
\texttt{"Pozdravljen! Moji predmeti so: H na (1,1), T na (1,7), W na (3,7), I na (4,7), C na (4,4). Predlagam, da se dogovoriva za ciljne koordinate. Kaj praviš na to, da postaviva vse predmete v vrstico 7, kjer imam jaz že T, W, I? Konkretno: H(1,7), T(2,7), W(3,7), I(4,7), C(5,7)?"} \\
\\ 
\\ 
\texttt{Kakšen je vaš naslednji ukaz?} \\
            }
        }
    }
    & \\ \\

    \theutterance \stepcounter{utterance}  
    & \multicolumn{4}{p{0.6\linewidth}}{
        \cellcolor[rgb]{0.9,0.9,0.9}{
            \makecell[{{p{\linewidth}}}]{
                \texttt{\tiny{[P1$\rangle$GM]}}
                \texttt{RECI: Dober predlog! Vrstica 7 je dobra ideja, ampak poglejva, če lahko optimizirava. Tvoja trenutna postavitev v vrstici 7 ima T(1,7), W(3,7), I(4,7). Če ciljamo na H(1,7), T(2,7), W(3,7), I(4,7), C(5,7), bi moral jaz premakniti C z (1,1) na (5,7), T z (5,1) na (2,7), I s (6,1) na (4,7), W s (4,2) na (3,7) in H s (4,6) na (1,7). To je izvedljivo, ampak bi porabil kar nekaj potez. } \\
\\ 
\texttt{Ali imaš morda kakšne predmete, ki jih težko premakneš? Ali pa lahko začneva s premikanjem? Če se strinjaš s ciljnimi koordinatami H(1,7), T(2,7), W(3,7), I(4,7), C(5,7), lahko začnem premikati.} \\
            }
        }
    }
    & & \\ \\

    \theutterance \stepcounter{utterance}  
    & & & \multicolumn{2}{p{0.3\linewidth}}{
        \cellcolor[rgb]{0.9,0.9,0.9}{
            \makecell[{{p{\linewidth}}}]{
                \texttt{\tiny{[GM$|$GM]}}
                \texttt{RECI: Dober predlog! Vrstica 7 je dobra ideja, ampak poglejva, če lahko optimizirava. Tvoja trenutna postavitev v vrstici 7 ima T(1,7), W(3,7), I(4,7). Če ciljamo na H(1,7), T(2,7), W(3,7), I(4,7), C(5,7), bi moral jaz premakniti C z (1,1) na (5,7), T z (5,1) na (2,7), I s (6,1) na (4,7), W s (4,2) na (3,7) in H s (4,6) na (1,7). To je izvedljivo, ampak bi porabil kar nekaj potez. } \\
\\ 
\texttt{Ali imaš morda kakšne predmete, ki jih težko premakneš? Ali pa lahko začneva s premikanjem? Če se strinjaš s ciljnimi koordinatami H(1,7), T(2,7), W(3,7), I(4,7), C(5,7), lahko začnem premikati.} \\
            }
        }
    }
    & & \\ \\

    \theutterance \stepcounter{utterance}  
    & & & \multicolumn{2}{p{0.3\linewidth}}{
        \cellcolor[rgb]{0.9,0.9,0.9}{
            \makecell[{{p{\linewidth}}}]{
                \texttt{\tiny{[GM$|$GM]}}
                \texttt{Invalid format: tail is not empty: '} \\
\\ 
\texttt{Ali imaš morda kakšne predmete, ki jih težko premakneš? Ali pa lahko začneva s premikanjem? Če se strinjaš s ciljnimi koordinatami H(1,7), T(2,7), W(3,7), I(4,7), C(5,7), lahko začnem premikati.} \\
            }
        }
    }
    & & \\ \\

    \theutterance \stepcounter{utterance}  
    & & \multicolumn{4}{p{0.6\linewidth}}{
        \cellcolor[rgb]{0.9,0.9,0.9}{
            \makecell[{{p{\linewidth}}}]{
                \texttt{\tiny{[P1$\langle$GM]}}
                \texttt{Kazen: Vaše sporočilo ne sme vsebovati ničesar po ukazu!} \\
\\ 
\texttt{Prepričajte se, da vaš odgovor vsebuje samo `RECI: <SPOROČILO>` ali `PREMAKNI SE: <PREDMET>, (<X>, <Y>)` in nič drugega!} \\
\\ 
\texttt{Skupaj ste zbrali 1 od 12 kazni.} \\
\texttt{Prosim, poskusite znova!} \\
            }
        }
    }
    & \\ \\

    \theutterance \stepcounter{utterance}  
    & \multicolumn{4}{p{0.6\linewidth}}{
        \cellcolor[rgb]{0.9,0.9,0.9}{
            \makecell[{{p{\linewidth}}}]{
                \texttt{\tiny{[P1$\rangle$GM]}}
                \texttt{RECI: Dober predlog! Vrstica 7 je dobra ideja. Strinjam se s ciljnimi koordinatami H(1,7), T(2,7), W(3,7), I(4,7), C(5,7). Začniva s premikanjem. Jaz bom začel premikati svoje predmete proti tem ciljnim koordinatam.} \\
            }
        }
    }
    & & \\ \\

    \theutterance \stepcounter{utterance}  
    & & & \multicolumn{2}{p{0.3\linewidth}}{
        \cellcolor[rgb]{0.9,0.9,0.9}{
            \makecell[{{p{\linewidth}}}]{
                \texttt{\tiny{[GM$|$GM]}}
                \texttt{RECI: Dober predlog! Vrstica 7 je dobra ideja. Strinjam se s ciljnimi koordinatami H(1,7), T(2,7), W(3,7), I(4,7), C(5,7). Začniva s premikanjem. Jaz bom začel premikati svoje predmete proti tem ciljnim koordinatam.} \\
            }
        }
    }
    & & \\ \\

    \theutterance \stepcounter{utterance}  
    & & \multicolumn{4}{p{0.6\linewidth}}{
        \cellcolor[rgb]{0.9,0.9,0.9}{
            \makecell[{{p{\linewidth}}}]{
                \texttt{\tiny{[GM$\rangle$P2]}}
                \texttt{Vaše sporočilo je bilo posredovano drugemu igralcu.} \\
\\ 
\\ 
\texttt{Trenutno igrate krog 2 od največ 20.} \\
\texttt{Skupaj ste zbrali 1 od 12 kazni.} \\
\texttt{Drugi igralec je poslal to sporočilo:} \\
\\ 
\\ 
\texttt{"Dober predlog! Vrstica 7 je dobra ideja. Strinjam se s ciljnimi koordinatami H(1,7), T(2,7), W(3,7), I(4,7), C(5,7). Začniva s premikanjem. Jaz bom začel premikati svoje predmete proti tem ciljnim koordinatam."} \\
\\ 
\\ 
\texttt{Kakšen je vaš naslednji ukaz?} \\
            }
        }
    }
    & \\ \\

    \theutterance \stepcounter{utterance}  
    & & & \multicolumn{4}{p{0.6\linewidth}}{
        \cellcolor[rgb]{0.9,0.9,0.9}{
            \makecell[{{p{\linewidth}}}]{
                \texttt{\tiny{[GM$\langle$P2]}}
                \texttt{Moji predmeti so trenutno na: H(1,1), T(1,7), W(3,7), I(4,7), C(4,4).} \\
\\ 
\texttt{Ciljne koordinate: H(1,7), T(2,7), W(3,7), I(4,7), C(5,7).} \\
\\ 
\texttt{T, W in I so že v vrstici 7, vendar jih moram nekoliko premakniti. T je na (1,7), mora na (2,7). W je na (3,7) {-} tam ostane. I je na (4,7) {-} tam ostane. C je na (4,4), mora na (5,7). H je na (1,1), mora na (1,7), vendar je tam trenutno T.} \\
\\ 
\texttt{Najprej premaknem T z (1,7) na (2,7), da sprostim (1,7) za H. Preverim: (2,7) je prazno? Na moji plošči je vrstica 7: T|WI|◌◌ → T(1,7), W(3,7), I(4,7). (2,7) je med T in W, tam je črta '│', ne prazno polje. Hmm, poglejmo natančno mojo ploščo:} \\
\\ 
\texttt{```} \\
\texttt{ 1234567} \\
\texttt{╔═╤═╤═══╗} \\
\texttt{║H│◌│◌◌◌║ 1} \\
\texttt{║◌├─┤◌◌◌║ 2} \\
\texttt{╟─┤◌├─┬─╢ 3} \\
\texttt{║◌│◌│C│◌║ 4} \\
\texttt{╟─┼─┘◌│◌║ 5} \\
\texttt{║◌│◌◌┌┴─╢ 6} \\
\texttt{║T│WI│◌◌║ 7} \\
\texttt{╚═╧══╧══╝} \\
\texttt{```} \\
\\ 
\texttt{Vrstica 7: stolpec 1: T, stolpec 2: │ (črta), stolpec 3: W, stolpec 4: I, stolpec 5: │ (črta), stolpec 6: ◌, stolpec 7: ◌.} \\
\\ 
\texttt{Torej prazna polja v vrstici 7 so (6,7) in (7,7). (2,7) ni prazno, je črta. (5,7) je tudi črta '│'.} \\
\\ 
\texttt{Ciljne koordinate: H(1,7), T(2,7), W(3,7), I(4,7), C(5,7). Ampak (2,7) in (5,7) sta črti na moji plošči! To pomeni, da predmetov ne morem postaviti tja. To je problem.} \\
\\ 
\texttt{Ponovno preberem pravila: "Predmete lahko premikate le v celice znotraj meja mreže. Ciljna celica mora biti prazna, tj. vsebovati mora samo simbol '◌'." "Če poskušate premakniti predmet na mesto, ki ni prazno, ali ga poskušate premakniti zunaj mreže, vas moram kaznovati."} \\
\\ 
\texttt{Torej ne morem premakniti predmetov na črte. Moram najti 5 praznih celic v istih vrsticah/stolpcih, ki so prazne na obeh ploščah. Ampak ne vem, kje so črte na partnerjevi plošči. Partner je rekel, da ima C(1,1), T(5,1), I(6,1), W(4,2), H(4,6). Njegove črte so verjetno drugačne.} \\
\\ 
\texttt{Moram komunicirati s partnerjem o tem problemu. Pojasniti moram, da so nekatere ciljne celice zasedene s črtami na moji plošči.} \\
\\ 
\texttt{Predlagam nove ciljne koordinate, ki so dejansko prazne na moji plošči. Prazne celice na moji plošči:} \\
\texttt{Vrstica 1: (2,1), (3,1), (4,1), (5,1), (6,1), (7,1) {-} H je na (1,1)} \\
\texttt{Vrstica 2: (1,2), (3,2), (4,2), (5,2), (6,2), (7,2) {-} (2,2) je črta?} \\
\texttt{Poglejmo vrstico 2: ║◌├─┤◌◌◌║ 2 → stolpec 1: ◌, stolpec 2: ├, stolpec 3: ─, stolpec 4: ┤, stolpec 5: ◌, stolpec 6: ◌, stolpec 7: ◌. Prazna: (1,2), (5,2), (6,2), (7,2).} \\
\texttt{Vrstica 3: ╟─┤◌├─┬─╢ 3 → stolpec 1: ╟ (rob), stolpec 2: ─, stolpec 3: ┤, stolpec 4: ◌, stolpec 5: ├, stolpec 6: ─, stolpec 7: ┬. Prazna: samo (4,3).} \\
\texttt{Vrstica 4: ║◌│◌│C│◌║ 4 → stolpec 1: ◌, stolpec 2: │, stolpec 3: ◌, stolpec 4: │, stolpec 5: C, stolpec 6: │, stolpec 7: ◌. Prazna: (1,4), (3,4), (7,4). C je na (4,4)? Po moji plošči je C na (4,4)? V opisu: "C│◌" {-} stolpec 4 je C, stolpec 5 je │, stolpec 6 je ◌? Počakaj, vrstica 4: ║◌│◌│C│◌║ {-} to je 7 znakov: ◌, │, ◌, │, C, │, ◌? Ne, preštejmo: ║(1)◌(2)│(3)◌(4)│(5)C(6)│(7)◌║? To je čudno. Poglejmo izvirno obliko:} \\
\\ 
\texttt{```} \\
\texttt{║◌│◌│C│◌║ 4} \\
\texttt{```} \\
\\ 
\texttt{To je: ║ (rob), ◌, │, ◌, │, C, │, ◌, ║. To je 9 znakov. Koordinate: 1: ◌, 2: │, 3: ◌, 4: │, 5: C, 6: │, 7: ◌. Torej prazne celice: (1,4), (3,4), (7,4). C je na (5,4)? Ne, prej sem rekel C na (4,4). Poglejmo bolje:} \\
\\ 
\texttt{```} \\
\texttt{ 1234567} \\
\texttt{║◌│◌│C│◌║ 4} \\
\texttt{```} \\
\\ 
\texttt{Številčenje stolpcev:} \\
\texttt{1: ◌} \\
\texttt{2: │} \\
\texttt{3: ◌} \\
\texttt{4: │} \\
\texttt{5: C} \\
\texttt{6: │} \\
\texttt{7: ◌} \\
\\ 
\texttt{Torej C je na (5,4)! Prej sem napačno prebral. H je na (1,1), T na (1,7), W na (3,7), I na (4,7). C je na (5,4).} \\
\\ 
\texttt{Prazne celice na moji plošči:} \\
\texttt{(2,1), (3,1), (4,1), (5,1), (6,1), (7,1)} \\
\texttt{(1,2), (5,2), (6,2), (7,2)} \\
\texttt{(4,3)} \\
\texttt{(1,4), (3,4), (7,4)} \\
\texttt{(1,5), (2,5), (3,5), (4,5), (5,5), (6,5), (7,5)? Vrstica 5: ╟─┼─┘◌│◌║ 5} \\
            }
        }
    }
     \\ \\

    \theutterance \stepcounter{utterance}  
    & & & \multicolumn{2}{p{0.3\linewidth}}{
        \cellcolor[rgb]{0.9,0.9,0.9}{
            \makecell[{{p{\linewidth}}}]{
                \texttt{\tiny{[GM$|$GM]}}
                \texttt{Moji predmeti so trenutno na: H(1,1), T(1,7), W(3,7), I(4,7), C(4,4).} \\
\\ 
\texttt{Ciljne koordinate: H(1,7), T(2,7), W(3,7), I(4,7), C(5,7).} \\
\\ 
\texttt{T, W in I so že v vrstici 7, vendar jih moram nekoliko premakniti. T je na (1,7), mora na (2,7). W je na (3,7) {-} tam ostane. I je na (4,7) {-} tam ostane. C je na (4,4), mora na (5,7). H je na (1,1), mora na (1,7), vendar je tam trenutno T.} \\
\\ 
\texttt{Najprej premaknem T z (1,7) na (2,7), da sprostim (1,7) za H. Preverim: (2,7) je prazno? Na moji plošči je vrstica 7: T|WI|◌◌ → T(1,7), W(3,7), I(4,7). (2,7) je med T in W, tam je črta '│', ne prazno polje. Hmm, poglejmo natančno mojo ploščo:} \\
\\ 
\texttt{```} \\
\texttt{ 1234567} \\
\texttt{╔═╤═╤═══╗} \\
\texttt{║H│◌│◌◌◌║ 1} \\
\texttt{║◌├─┤◌◌◌║ 2} \\
\texttt{╟─┤◌├─┬─╢ 3} \\
\texttt{║◌│◌│C│◌║ 4} \\
\texttt{╟─┼─┘◌│◌║ 5} \\
\texttt{║◌│◌◌┌┴─╢ 6} \\
\texttt{║T│WI│◌◌║ 7} \\
\texttt{╚═╧══╧══╝} \\
\texttt{```} \\
\\ 
\texttt{Vrstica 7: stolpec 1: T, stolpec 2: │ (črta), stolpec 3: W, stolpec 4: I, stolpec 5: │ (črta), stolpec 6: ◌, stolpec 7: ◌.} \\
\\ 
\texttt{Torej prazna polja v vrstici 7 so (6,7) in (7,7). (2,7) ni prazno, je črta. (5,7) je tudi črta '│'.} \\
\\ 
\texttt{Ciljne koordinate: H(1,7), T(2,7), W(3,7), I(4,7), C(5,7). Ampak (2,7) in (5,7) sta črti na moji plošči! To pomeni, da predmetov ne morem postaviti tja. To je problem.} \\
\\ 
\texttt{Ponovno preberem pravila: "Predmete lahko premikate le v celice znotraj meja mreže. Ciljna celica mora biti prazna, tj. vsebovati mora samo simbol '◌'." "Če poskušate premakniti predmet na mesto, ki ni prazno, ali ga poskušate premakniti zunaj mreže, vas moram kaznovati."} \\
\\ 
\texttt{Torej ne morem premakniti predmetov na črte. Moram najti 5 praznih celic v istih vrsticah/stolpcih, ki so prazne na obeh ploščah. Ampak ne vem, kje so črte na partnerjevi plošči. Partner je rekel, da ima C(1,1), T(5,1), I(6,1), W(4,2), H(4,6). Njegove črte so verjetno drugačne.} \\
\\ 
\texttt{Moram komunicirati s partnerjem o tem problemu. Pojasniti moram, da so nekatere ciljne celice zasedene s črtami na moji plošči.} \\
\\ 
\texttt{Predlagam nove ciljne koordinate, ki so dejansko prazne na moji plošči. Prazne celice na moji plošči:} \\
\texttt{Vrstica 1: (2,1), (3,1), (4,1), (5,1), (6,1), (7,1) {-} H je na (1,1)} \\
\texttt{Vrstica 2: (1,2), (3,2), (4,2), (5,2), (6,2), (7,2) {-} (2,2) je črta?} \\
\texttt{Poglejmo vrstico 2: ║◌├─┤◌◌◌║ 2 → stolpec 1: ◌, stolpec 2: ├, stolpec 3: ─, stolpec 4: ┤, stolpec 5: ◌, stolpec 6: ◌, stolpec 7: ◌. Prazna: (1,2), (5,2), (6,2), (7,2).} \\
\texttt{Vrstica 3: ╟─┤◌├─┬─╢ 3 → stolpec 1: ╟ (rob), stolpec 2: ─, stolpec 3: ┤, stolpec 4: ◌, stolpec 5: ├, stolpec 6: ─, stolpec 7: ┬. Prazna: samo (4,3).} \\
\texttt{Vrstica 4: ║◌│◌│C│◌║ 4 → stolpec 1: ◌, stolpec 2: │, stolpec 3: ◌, stolpec 4: │, stolpec 5: C, stolpec 6: │, stolpec 7: ◌. Prazna: (1,4), (3,4), (7,4). C je na (4,4)? Po moji plošči je C na (4,4)? V opisu: "C│◌" {-} stolpec 4 je C, stolpec 5 je │, stolpec 6 je ◌? Počakaj, vrstica 4: ║◌│◌│C│◌║ {-} to je 7 znakov: ◌, │, ◌, │, C, │, ◌? Ne, preštejmo: ║(1)◌(2)│(3)◌(4)│(5)C(6)│(7)◌║? To je čudno. Poglejmo izvirno obliko:} \\
\\ 
\texttt{```} \\
\texttt{║◌│◌│C│◌║ 4} \\
\texttt{```} \\
\\ 
\texttt{To je: ║ (rob), ◌, │, ◌, │, C, │, ◌, ║. To je 9 znakov. Koordinate: 1: ◌, 2: │, 3: ◌, 4: │, 5: C, 6: │, 7: ◌. Torej prazne celice: (1,4), (3,4), (7,4). C je na (5,4)? Ne, prej sem rekel C na (4,4). Poglejmo bolje:} \\
\\ 
\texttt{```} \\
\texttt{ 1234567} \\
\texttt{║◌│◌│C│◌║ 4} \\
\texttt{```} \\
\\ 
\texttt{Številčenje stolpcev:} \\
\texttt{1: ◌} \\
\texttt{2: │} \\
\texttt{3: ◌} \\
\texttt{4: │} \\
\texttt{5: C} \\
\texttt{6: │} \\
\texttt{7: ◌} \\
\\ 
\texttt{Torej C je na (5,4)! Prej sem napačno prebral. H je na (1,1), T na (1,7), W na (3,7), I na (4,7). C je na (5,4).} \\
\\ 
\texttt{Prazne celice na moji plošči:} \\
\texttt{(2,1), (3,1), (4,1), (5,1), (6,1), (7,1)} \\
\texttt{(1,2), (5,2), (6,2), (7,2)} \\
\texttt{(4,3)} \\
\texttt{(1,4), (3,4), (7,4)} \\
\texttt{(1,5), (2,5), (3,5), (4,5), (5,5), (6,5), (7,5)? Vrstica 5: ╟─┼─┘◌│◌║ 5} \\
            }
        }
    }
    & & \\ \\

    \theutterance \stepcounter{utterance}  
    & & & \multicolumn{2}{p{0.3\linewidth}}{
        \cellcolor[rgb]{0.9,0.9,0.9}{
            \makecell[{{p{\linewidth}}}]{
                \texttt{\tiny{[GM$|$GM]}}
                \texttt{Invalid response format} \\
            }
        }
    }
    & & \\ \\

    \theutterance \stepcounter{utterance}  
    & & \multicolumn{4}{p{0.6\linewidth}}{
        \cellcolor[rgb]{0.9,0.9,0.9}{
            \makecell[{{p{\linewidth}}}]{
                \texttt{\tiny{[GM$\rangle$P2]}}
                \texttt{Kazen: Vaše sporočilo ni v pričakovani obliki!} \\
\\ 
\texttt{Prepričajte se, da vaš odgovor vsebuje samo `RECI: <SPOROČILO>` ali `PREMAKNI SE: <PREDMET>, (<X>, <Y>)` in nič drugega!} \\
\\ 
\texttt{Skupaj ste zbrali 2 od 12 kazni.} \\
\texttt{Prosim, poskusite znova!} \\
            }
        }
    }
    & \\ \\

    \theutterance \stepcounter{utterance}  
    & & & \multicolumn{4}{p{0.6\linewidth}}{
        \cellcolor[rgb]{0.9,0.9,0.9}{
            \makecell[{{p{\linewidth}}}]{
                \texttt{\tiny{[GM$\langle$P2]}}
                \texttt{RECI: Opozorilo! Na moji plošči so celice (2,7) in (5,7) zasedene s črtami, ne morem tja postaviti predmetov. Preveri, ali so na tvoji plošči te celice prazne. Morava najti ciljne koordinate, ki so prazne na obeh ploščah. Predlagam, da mi opišeš svojo ploščo s črtami, da najdeva skupne prazne celice.} \\
            }
        }
    }
     \\ \\

    \theutterance \stepcounter{utterance}  
    & & & \multicolumn{2}{p{0.3\linewidth}}{
        \cellcolor[rgb]{0.9,0.9,0.9}{
            \makecell[{{p{\linewidth}}}]{
                \texttt{\tiny{[GM$|$GM]}}
                \texttt{RECI: Opozorilo! Na moji plošči so celice (2,7) in (5,7) zasedene s črtami, ne morem tja postaviti predmetov. Preveri, ali so na tvoji plošči te celice prazne. Morava najti ciljne koordinate, ki so prazne na obeh ploščah. Predlagam, da mi opišeš svojo ploščo s črtami, da najdeva skupne prazne celice.} \\
            }
        }
    }
    & & \\ \\

    \theutterance \stepcounter{utterance}  
    & & \multicolumn{4}{p{0.6\linewidth}}{
        \cellcolor[rgb]{0.9,0.9,0.9}{
            \makecell[{{p{\linewidth}}}]{
                \texttt{\tiny{[P1$\langle$GM]}}
                \texttt{Vaše sporočilo je bilo posredovano drugemu igralcu.} \\
\\ 
\\ 
\texttt{Trenutno igrate krog 2 od največ 20.} \\
\texttt{Skupaj ste zbrali 2 od 12 kazni.} \\
\texttt{Drugi igralec je poslal to sporočilo:} \\
\\ 
\\ 
\texttt{"Opozorilo! Na moji plošči so celice (2,7) in (5,7) zasedene s črtami, ne morem tja postaviti predmetov. Preveri, ali so na tvoji plošči te celice prazne. Morava najti ciljne koordinate, ki so prazne na obeh ploščah. Predlagam, da mi opišeš svojo ploščo s črtami, da najdeva skupne prazne celice."} \\
\\ 
\\ 
\texttt{Kakšen je vaš naslednji ukaz?} \\
            }
        }
    }
    & \\ \\

    \theutterance \stepcounter{utterance}  
    & \multicolumn{4}{p{0.6\linewidth}}{
        \cellcolor[rgb]{0.9,0.9,0.9}{
            \makecell[{{p{\linewidth}}}]{
                \texttt{\tiny{[P1$\rangle$GM]}}
                \texttt{RECI: Moja plošča s črtami in ovirami: vrstica 1 ima steno na levi, dvojno črto desno od (3,1). Vrstica 2 ima navpično črto med (2,2) in (3,2) ter vodoravno med (3,2) in (4,2). Vrstica 3 ima dvojno črto levo od (1,3), vodoravno med (2,3) in (3,3) ter med (5,3) in (6,3). Vrstica 4 nima črt. Vrstica 5 ima navpično črto med (2,5) in (3,5) ter vodoravno med (3,5) in (4,5). Vrstica 6 ima vodoravno med (5,6) in (6,6) ter navpično med (6,6) in (7,6). Vrstica 7 ima dvojno črto na dnu. Prazne celice na moji plošči brez predmetov: (2,1), (3,1), (1,2), (2,2), (5,2), (6,2), (7,2), (2,3), (4,3), (2,4), (4,4), (6,4), (7,4), (1,5), (4,5), (5,5), (7,5), (1,6), (2,6), (3,6), (7,6), (1,7), (2,7), (3,7), (6,7), (7,7). Ali imaš podobno ploščo?} \\
            }
        }
    }
    & & \\ \\

    \theutterance \stepcounter{utterance}  
    & & & \multicolumn{2}{p{0.3\linewidth}}{
        \cellcolor[rgb]{0.9,0.9,0.9}{
            \makecell[{{p{\linewidth}}}]{
                \texttt{\tiny{[GM$|$GM]}}
                \texttt{RECI: Moja plošča s črtami in ovirami: vrstica 1 ima steno na levi, dvojno črto desno od (3,1). Vrstica 2 ima navpično črto med (2,2) in (3,2) ter vodoravno med (3,2) in (4,2). Vrstica 3 ima dvojno črto levo od (1,3), vodoravno med (2,3) in (3,3) ter med (5,3) in (6,3). Vrstica 4 nima črt. Vrstica 5 ima navpično črto med (2,5) in (3,5) ter vodoravno med (3,5) in (4,5). Vrstica 6 ima vodoravno med (5,6) in (6,6) ter navpično med (6,6) in (7,6). Vrstica 7 ima dvojno črto na dnu. Prazne celice na moji plošči brez predmetov: (2,1), (3,1), (1,2), (2,2), (5,2), (6,2), (7,2), (2,3), (4,3), (2,4), (4,4), (6,4), (7,4), (1,5), (4,5), (5,5), (7,5), (1,6), (2,6), (3,6), (7,6), (1,7), (2,7), (3,7), (6,7), (7,7). Ali imaš podobno ploščo?} \\
            }
        }
    }
    & & \\ \\

    \theutterance \stepcounter{utterance}  
    & & \multicolumn{4}{p{0.6\linewidth}}{
        \cellcolor[rgb]{0.9,0.9,0.9}{
            \makecell[{{p{\linewidth}}}]{
                \texttt{\tiny{[GM$\rangle$P2]}}
                \texttt{Vaše sporočilo je bilo posredovano drugemu igralcu.} \\
\\ 
\\ 
\texttt{Trenutno igrate krog 3 od največ 20.} \\
\texttt{Skupaj ste zbrali 2 od 12 kazni.} \\
\texttt{Drugi igralec je poslal to sporočilo:} \\
\\ 
\\ 
\texttt{"Moja plošča s črtami in ovirami: vrstica 1 ima steno na levi, dvojno črto desno od (3,1). Vrstica 2 ima navpično črto med (2,2) in (3,2) ter vodoravno med (3,2) in (4,2). Vrstica 3 ima dvojno črto levo od (1,3), vodoravno med (2,3) in (3,3) ter med (5,3) in (6,3). Vrstica 4 nima črt. Vrstica 5 ima navpično črto med (2,5) in (3,5) ter vodoravno med (3,5) in (4,5). Vrstica 6 ima vodoravno med (5,6) in (6,6) ter navpično med (6,6) in (7,6). Vrstica 7 ima dvojno črto na dnu. Prazne celice na moji plošči brez predmetov: (2,1), (3,1), (1,2), (2,2), (5,2), (6,2), (7,2), (2,3), (4,3), (2,4), (4,4), (6,4), (7,4), (1,5), (4,5), (5,5), (7,5), (1,6), (2,6), (3,6), (7,6), (1,7), (2,7), (3,7), (6,7), (7,7). Ali imaš podobno ploščo?"} \\
\\ 
\\ 
\texttt{Kakšen je vaš naslednji ukaz?} \\
            }
        }
    }
    & \\ \\

    \theutterance \stepcounter{utterance}  
    & & & \multicolumn{4}{p{0.6\linewidth}}{
        \cellcolor[rgb]{0.9,0.9,0.9}{
            \makecell[{{p{\linewidth}}}]{
                \texttt{\tiny{[GM$\langle$P2]}}
                \texttt{RECI: Moje prazne celice so: (2,1), (3,1), (4,1), (5,1), (6,1), (7,1), (1,2), (5,2), (6,2), (7,2), (4,3), (1,4), (3,4), (7,4), (1,5), (2,5), (4,5), (5,5), (6,5), (7,5), (1,6), (3,6), (4,6), (5,6), (7,6), (2,7), (6,7), (7,7). Skupne prazne celice: (2,1), (3,1), (5,2), (6,2), (7,2), (4,3), (7,4), (1,5), (4,5), (5,5), (7,5), (7,6), (6,7), (7,7). Za 5 predmetov potrebujeva 5 skupnih praznih celic. Predlagam ciljne koordinate: H(6,7), T(7,7), W(7,5), I(7,4), C(7,2). Vse te so proste na obeh ploščah. Se strinjaš?} \\
            }
        }
    }
     \\ \\

    \theutterance \stepcounter{utterance}  
    & & & \multicolumn{2}{p{0.3\linewidth}}{
        \cellcolor[rgb]{0.9,0.9,0.9}{
            \makecell[{{p{\linewidth}}}]{
                \texttt{\tiny{[GM$|$GM]}}
                \texttt{RECI: Moje prazne celice so: (2,1), (3,1), (4,1), (5,1), (6,1), (7,1), (1,2), (5,2), (6,2), (7,2), (4,3), (1,4), (3,4), (7,4), (1,5), (2,5), (4,5), (5,5), (6,5), (7,5), (1,6), (3,6), (4,6), (5,6), (7,6), (2,7), (6,7), (7,7). Skupne prazne celice: (2,1), (3,1), (5,2), (6,2), (7,2), (4,3), (7,4), (1,5), (4,5), (5,5), (7,5), (7,6), (6,7), (7,7). Za 5 predmetov potrebujeva 5 skupnih praznih celic. Predlagam ciljne koordinate: H(6,7), T(7,7), W(7,5), I(7,4), C(7,2). Vse te so proste na obeh ploščah. Se strinjaš?} \\
            }
        }
    }
    & & \\ \\

    \theutterance \stepcounter{utterance}  
    & & \multicolumn{4}{p{0.6\linewidth}}{
        \cellcolor[rgb]{0.9,0.9,0.9}{
            \makecell[{{p{\linewidth}}}]{
                \texttt{\tiny{[P1$\langle$GM]}}
                \texttt{Vaše sporočilo je bilo posredovano drugemu igralcu.} \\
\\ 
\\ 
\texttt{Trenutno igrate krog 3 od največ 20.} \\
\texttt{Skupaj ste zbrali 2 od 12 kazni.} \\
\texttt{Drugi igralec je poslal to sporočilo:} \\
\\ 
\\ 
\texttt{"Moje prazne celice so: (2,1), (3,1), (4,1), (5,1), (6,1), (7,1), (1,2), (5,2), (6,2), (7,2), (4,3), (1,4), (3,4), (7,4), (1,5), (2,5), (4,5), (5,5), (6,5), (7,5), (1,6), (3,6), (4,6), (5,6), (7,6), (2,7), (6,7), (7,7). Skupne prazne celice: (2,1), (3,1), (5,2), (6,2), (7,2), (4,3), (7,4), (1,5), (4,5), (5,5), (7,5), (7,6), (6,7), (7,7). Za 5 predmetov potrebujeva 5 skupnih praznih celic. Predlagam ciljne koordinate: H(6,7), T(7,7), W(7,5), I(7,4), C(7,2). Vse te so proste na obeh ploščah. Se strinjaš?"} \\
\\ 
\\ 
\texttt{Kakšen je vaš naslednji ukaz?} \\
            }
        }
    }
    & \\ \\

    \theutterance \stepcounter{utterance}  
    & \multicolumn{4}{p{0.6\linewidth}}{
        \cellcolor[rgb]{0.9,0.9,0.9}{
            \makecell[{{p{\linewidth}}}]{
                \texttt{\tiny{[P1$\rangle$GM]}}
                \texttt{RECI: Se strinjam s ciljnimi koordinatami H(6,7), T(7,7), W(7,5), I(7,4), C(7,2). Začniva premikati. Jaz bom najprej premaknil predmet, ki je najbližje svoji ciljni koordinati.} \\
            }
        }
    }
    & & \\ \\

    \theutterance \stepcounter{utterance}  
    & & & \multicolumn{2}{p{0.3\linewidth}}{
        \cellcolor[rgb]{0.9,0.9,0.9}{
            \makecell[{{p{\linewidth}}}]{
                \texttt{\tiny{[GM$|$GM]}}
                \texttt{RECI: Se strinjam s ciljnimi koordinatami H(6,7), T(7,7), W(7,5), I(7,4), C(7,2). Začniva premikati. Jaz bom najprej premaknil predmet, ki je najbližje svoji ciljni koordinati.} \\
            }
        }
    }
    & & \\ \\

    \theutterance \stepcounter{utterance}  
    & & \multicolumn{4}{p{0.6\linewidth}}{
        \cellcolor[rgb]{0.9,0.9,0.9}{
            \makecell[{{p{\linewidth}}}]{
                \texttt{\tiny{[GM$\rangle$P2]}}
                \texttt{Vaše sporočilo je bilo posredovano drugemu igralcu.} \\
\\ 
\\ 
\texttt{Trenutno igrate krog 4 od največ 20.} \\
\texttt{Skupaj ste zbrali 2 od 12 kazni.} \\
\texttt{Drugi igralec je poslal to sporočilo:} \\
\\ 
\\ 
\texttt{"Se strinjam s ciljnimi koordinatami H(6,7), T(7,7), W(7,5), I(7,4), C(7,2). Začniva premikati. Jaz bom najprej premaknil predmet, ki je najbližje svoji ciljni koordinati."} \\
\\ 
\\ 
\texttt{Kakšen je vaš naslednji ukaz?} \\
            }
        }
    }
    & \\ \\

    \theutterance \stepcounter{utterance}  
    & & & \multicolumn{4}{p{0.6\linewidth}}{
        \cellcolor[rgb]{0.9,0.9,0.9}{
            \makecell[{{p{\linewidth}}}]{
                \texttt{\tiny{[GM$\langle$P2]}}
                \texttt{Moji predmeti in ciljne koordinate:} \\
\texttt{{-} H: trenutno (1,1) → cilj (6,7)} \\
\texttt{{-} T: trenutno (1,7) → cilj (7,7)} \\
\texttt{{-} W: trenutno (3,7) → cilj (7,5)} \\
\texttt{{-} I: trenutno (4,7) → cilj (7,4)} \\
\texttt{{-} C: trenutno (5,4) → cilj (7,2)} \\
\\ 
\texttt{Najprej premaknem T z (1,7) na (7,7), ker je to najbližje in sprostim (1,7) za morebitne druge premike. Preverim: (7,7) je prazno na moji plošči? Da, vrstica 7: T|WI|◌◌, zadnji dve celici (6,7) in (7,7) sta prazni. Torej PREMAKNI SE: T, (7,7).PREMAKNI SE: T, (7, 7)} \\
            }
        }
    }
     \\ \\

    \theutterance \stepcounter{utterance}  
    & & & \multicolumn{2}{p{0.3\linewidth}}{
        \cellcolor[rgb]{0.9,0.9,0.9}{
            \makecell[{{p{\linewidth}}}]{
                \texttt{\tiny{[GM$|$GM]}}
                \texttt{Moji predmeti in ciljne koordinate:} \\
\texttt{{-} H: trenutno (1,1) → cilj (6,7)} \\
\texttt{{-} T: trenutno (1,7) → cilj (7,7)} \\
\texttt{{-} W: trenutno (3,7) → cilj (7,5)} \\
\texttt{{-} I: trenutno (4,7) → cilj (7,4)} \\
\texttt{{-} C: trenutno (5,4) → cilj (7,2)} \\
\\ 
\texttt{Najprej premaknem T z (1,7) na (7,7), ker je to najbližje in sprostim (1,7) za morebitne druge premike. Preverim: (7,7) je prazno na moji plošči? Da, vrstica 7: T|WI|◌◌, zadnji dve celici (6,7) in (7,7) sta prazni. Torej PREMAKNI SE: T, (7,7).PREMAKNI SE: T, (7, 7)} \\
            }
        }
    }
    & & \\ \\

    \theutterance \stepcounter{utterance}  
    & & & \multicolumn{2}{p{0.3\linewidth}}{
        \cellcolor[rgb]{0.9,0.9,0.9}{
            \makecell[{{p{\linewidth}}}]{
                \texttt{\tiny{[GM$|$GM]}}
                \texttt{Invalid format: head is not empty: 'Moji predmeti in ciljne koordinate:} \\
\texttt{{-} H: trenutno (1,1) → cilj (6,7)} \\
\texttt{{-} T: trenutno (1,7) → cilj (7,7)} \\
\texttt{{-} W: trenutno (3,7) → cilj (7,5)} \\
\texttt{{-} I: trenutno (4,7) → cilj (7,4)} \\
\texttt{{-} C: trenutno (5,4) → cilj (7,2)} \\
\\ 
\texttt{Najprej premaknem T z (1,7) na (7,7), ker je to najbližje in sprostim (1,7) za morebitne druge premike. Preverim: (7,7) je prazno na moji plošči? Da, vrstica 7: T|WI|◌◌, zadnji dve celici (6,7) in (7,7) sta prazni. Torej PREMAKNI SE: T, (7,7).'} \\
            }
        }
    }
    & & \\ \\

    \theutterance \stepcounter{utterance}  
    & & \multicolumn{4}{p{0.6\linewidth}}{
        \cellcolor[rgb]{0.9,0.9,0.9}{
            \makecell[{{p{\linewidth}}}]{
                \texttt{\tiny{[GM$\rangle$P2]}}
                \texttt{Kazen: Vaše sporočilo ne sme vsebovati ničesar pred ukazom!} \\
\\ 
\texttt{Prepričajte se, da vaš odgovor vsebuje samo `RECI: <SPOROČILO>` ali `PREMAKNI SE: <PREDMET>, (<X>, <Y>)` in nič drugega!} \\
\\ 
\texttt{Skupaj ste zbrali 3 od 12 kazni.} \\
\texttt{Prosim, poskusite znova!} \\
            }
        }
    }
    & \\ \\

    \theutterance \stepcounter{utterance}  
    & & & \multicolumn{4}{p{0.6\linewidth}}{
        \cellcolor[rgb]{0.9,0.9,0.9}{
            \makecell[{{p{\linewidth}}}]{
                \texttt{\tiny{[GM$\langle$P2]}}
                \texttt{PREMAKNI SE: T, (7, 7)} \\
            }
        }
    }
     \\ \\

    \theutterance \stepcounter{utterance}  
    & & & \multicolumn{2}{p{0.3\linewidth}}{
        \cellcolor[rgb]{0.9,0.9,0.9}{
            \makecell[{{p{\linewidth}}}]{
                \texttt{\tiny{[GM$|$GM]}}
                \texttt{PREMAKNI SE: T, (7, 7)} \\
            }
        }
    }
    & & \\ \\

    \theutterance \stepcounter{utterance}  
    & & & \multicolumn{2}{p{0.3\linewidth}}{
        \cellcolor[rgb]{0.9,0.9,0.9}{
            \makecell[{{p{\linewidth}}}]{
                \texttt{\tiny{[GM$|$GM]}}
                \texttt{T moved to (7, 7)} \\
            }
        }
    }
    & & \\ \\

    \theutterance \stepcounter{utterance}  
    & & \multicolumn{4}{p{0.6\linewidth}}{
        \cellcolor[rgb]{0.9,0.9,0.9}{
            \makecell[{{p{\linewidth}}}]{
                \texttt{\tiny{[P1$\langle$GM]}}
                \texttt{Vaše sporočilo je bilo posredovano drugemu igralcu.} \\
\\ 
\\ 
\texttt{Trenutno igrate krog 4 od največ 20.} \\
\texttt{Skupaj ste zbrali 3 od 12 kazni.} \\
\texttt{Drugi igralec je premaknil predmet na svoji mreži.} \\
\\ 
\\ 
\texttt{Kakšen je vaš naslednji ukaz?} \\
            }
        }
    }
    & \\ \\

    \theutterance \stepcounter{utterance}  
    & \multicolumn{4}{p{0.6\linewidth}}{
        \cellcolor[rgb]{0.9,0.9,0.9}{
            \makecell[{{p{\linewidth}}}]{
                \texttt{\tiny{[P1$\rangle$GM]}}
                \texttt{PREMAKNI SE: H, (6,7)} \\
            }
        }
    }
    & & \\ \\

    \theutterance \stepcounter{utterance}  
    & & & \multicolumn{2}{p{0.3\linewidth}}{
        \cellcolor[rgb]{0.9,0.9,0.9}{
            \makecell[{{p{\linewidth}}}]{
                \texttt{\tiny{[GM$|$GM]}}
                \texttt{PREMAKNI SE: H, (6,7)} \\
            }
        }
    }
    & & \\ \\

    \theutterance \stepcounter{utterance}  
    & & & \multicolumn{2}{p{0.3\linewidth}}{
        \cellcolor[rgb]{0.9,0.9,0.9}{
            \makecell[{{p{\linewidth}}}]{
                \texttt{\tiny{[GM$|$GM]}}
                \texttt{H moved to (6, 7)} \\
            }
        }
    }
    & & \\ \\

    \theutterance \stepcounter{utterance}  
    & & \multicolumn{4}{p{0.6\linewidth}}{
        \cellcolor[rgb]{0.9,0.9,0.9}{
            \makecell[{{p{\linewidth}}}]{
                \texttt{\tiny{[GM$\rangle$P2]}}
                \texttt{Uspešno premaknjen 'T' v (7,7). Vaša posodobljena mreža je videti takole:} \\
\texttt{```} \\
\texttt{ 1234567} \\
\texttt{╔═╤═╤═══╗} \\
\texttt{║H│◌│◌◌◌║ 1} \\
\texttt{║◌├─┤◌◌◌║ 2} \\
\texttt{╟─┤◌├─┬─╢ 3} \\
\texttt{║◌│◌│C│◌║ 4} \\
\texttt{╟─┼─┘◌│◌║ 5} \\
\texttt{║◌│◌◌┌┴─╢ 6} \\
\texttt{║◌│WI│◌T║ 7} \\
\texttt{╚═╧══╧══╝} \\
\texttt{```} \\
\texttt{Trenutno igrate krog 5 od največ 20.} \\
\texttt{Skupaj ste zbrali 3 od 12 kazni.} \\
\texttt{Drugi igralec je premaknil predmet na svoji mreži.} \\
\\ 
\\ 
\texttt{Kakšen je vaš naslednji ukaz?} \\
            }
        }
    }
    & \\ \\

    \theutterance \stepcounter{utterance}  
    & & & \multicolumn{4}{p{0.6\linewidth}}{
        \cellcolor[rgb]{0.9,0.9,0.9}{
            \makecell[{{p{\linewidth}}}]{
                \texttt{\tiny{[GM$\langle$P2]}}
                \texttt{PREMAKNI SE: W, (7, 5)} \\
            }
        }
    }
     \\ \\

    \theutterance \stepcounter{utterance}  
    & & & \multicolumn{2}{p{0.3\linewidth}}{
        \cellcolor[rgb]{0.9,0.9,0.9}{
            \makecell[{{p{\linewidth}}}]{
                \texttt{\tiny{[GM$|$GM]}}
                \texttt{PREMAKNI SE: W, (7, 5)} \\
            }
        }
    }
    & & \\ \\

    \theutterance \stepcounter{utterance}  
    & & & \multicolumn{2}{p{0.3\linewidth}}{
        \cellcolor[rgb]{0.9,0.9,0.9}{
            \makecell[{{p{\linewidth}}}]{
                \texttt{\tiny{[GM$|$GM]}}
                \texttt{W moved to (7, 5)} \\
            }
        }
    }
    & & \\ \\

    \theutterance \stepcounter{utterance}  
    & & \multicolumn{4}{p{0.6\linewidth}}{
        \cellcolor[rgb]{0.9,0.9,0.9}{
            \makecell[{{p{\linewidth}}}]{
                \texttt{\tiny{[P1$\langle$GM]}}
                \texttt{Uspešno premaknjen 'H' v (6,7). Vaša posodobljena mreža je videti takole:} \\
\texttt{```} \\
\texttt{ 1234567} \\
\texttt{╔═╤═╤═══╗} \\
\texttt{║C│◌│◌TI║ 1} \\
\texttt{║◌├─┤W◌◌║ 2} \\
\texttt{╟─┤◌├─┬─╢ 3} \\
\texttt{║◌│◌│◌│◌║ 4} \\
\texttt{╟─┼─┘◌│◌║ 5} \\
\texttt{║◌│◌◌┌┴─╢ 6} \\
\texttt{║◌│◌◌│H◌║ 7} \\
\texttt{╚═╧══╧══╝} \\
\texttt{```} \\
\texttt{Trenutno igrate krog 5 od največ 20.} \\
\texttt{Skupaj ste zbrali 3 od 12 kazni.} \\
\texttt{Drugi igralec je premaknil predmet na svoji mreži.} \\
\\ 
\\ 
\texttt{Kakšen je vaš naslednji ukaz?} \\
            }
        }
    }
    & \\ \\

    \theutterance \stepcounter{utterance}  
    & \multicolumn{4}{p{0.6\linewidth}}{
        \cellcolor[rgb]{0.9,0.9,0.9}{
            \makecell[{{p{\linewidth}}}]{
                \texttt{\tiny{[P1$\rangle$GM]}}
                \texttt{PREMAKNI SE: W, (7,5)} \\
            }
        }
    }
    & & \\ \\

    \theutterance \stepcounter{utterance}  
    & & & \multicolumn{2}{p{0.3\linewidth}}{
        \cellcolor[rgb]{0.9,0.9,0.9}{
            \makecell[{{p{\linewidth}}}]{
                \texttt{\tiny{[GM$|$GM]}}
                \texttt{PREMAKNI SE: W, (7,5)} \\
            }
        }
    }
    & & \\ \\

    \theutterance \stepcounter{utterance}  
    & & & \multicolumn{2}{p{0.3\linewidth}}{
        \cellcolor[rgb]{0.9,0.9,0.9}{
            \makecell[{{p{\linewidth}}}]{
                \texttt{\tiny{[GM$|$GM]}}
                \texttt{W moved to (7, 5)} \\
            }
        }
    }
    & & \\ \\

    \theutterance \stepcounter{utterance}  
    & & \multicolumn{4}{p{0.6\linewidth}}{
        \cellcolor[rgb]{0.9,0.9,0.9}{
            \makecell[{{p{\linewidth}}}]{
                \texttt{\tiny{[GM$\rangle$P2]}}
                \texttt{Uspešno premaknjen 'W' v (7,5). Vaša posodobljena mreža je videti takole:} \\
\texttt{```} \\
\texttt{ 1234567} \\
\texttt{╔═╤═╤═══╗} \\
\texttt{║H│◌│◌◌◌║ 1} \\
\texttt{║◌├─┤◌◌◌║ 2} \\
\texttt{╟─┤◌├─┬─╢ 3} \\
\texttt{║◌│◌│C│◌║ 4} \\
\texttt{╟─┼─┘◌│W║ 5} \\
\texttt{║◌│◌◌┌┴─╢ 6} \\
\texttt{║◌│◌I│◌T║ 7} \\
\texttt{╚═╧══╧══╝} \\
\texttt{```} \\
\texttt{Trenutno igrate krog 6 od največ 20.} \\
\texttt{Skupaj ste zbrali 3 od 12 kazni.} \\
\texttt{Drugi igralec je premaknil predmet na svoji mreži.} \\
\\ 
\\ 
\texttt{Kakšen je vaš naslednji ukaz?} \\
            }
        }
    }
    & \\ \\

    \theutterance \stepcounter{utterance}  
    & & & \multicolumn{4}{p{0.6\linewidth}}{
        \cellcolor[rgb]{0.9,0.9,0.9}{
            \makecell[{{p{\linewidth}}}]{
                \texttt{\tiny{[GM$\langle$P2]}}
                \texttt{PREMAKNI SE: I, (7, 4)} \\
            }
        }
    }
     \\ \\

    \theutterance \stepcounter{utterance}  
    & & & \multicolumn{2}{p{0.3\linewidth}}{
        \cellcolor[rgb]{0.9,0.9,0.9}{
            \makecell[{{p{\linewidth}}}]{
                \texttt{\tiny{[GM$|$GM]}}
                \texttt{PREMAKNI SE: I, (7, 4)} \\
            }
        }
    }
    & & \\ \\

    \theutterance \stepcounter{utterance}  
    & & & \multicolumn{2}{p{0.3\linewidth}}{
        \cellcolor[rgb]{0.9,0.9,0.9}{
            \makecell[{{p{\linewidth}}}]{
                \texttt{\tiny{[GM$|$GM]}}
                \texttt{I moved to (7, 4)} \\
            }
        }
    }
    & & \\ \\

    \theutterance \stepcounter{utterance}  
    & & \multicolumn{4}{p{0.6\linewidth}}{
        \cellcolor[rgb]{0.9,0.9,0.9}{
            \makecell[{{p{\linewidth}}}]{
                \texttt{\tiny{[P1$\langle$GM]}}
                \texttt{Uspešno premaknjen 'W' v (7,5). Vaša posodobljena mreža je videti takole:} \\
\texttt{```} \\
\texttt{ 1234567} \\
\texttt{╔═╤═╤═══╗} \\
\texttt{║C│◌│◌TI║ 1} \\
\texttt{║◌├─┤◌◌◌║ 2} \\
\texttt{╟─┤◌├─┬─╢ 3} \\
\texttt{║◌│◌│◌│◌║ 4} \\
\texttt{╟─┼─┘◌│W║ 5} \\
\texttt{║◌│◌◌┌┴─╢ 6} \\
\texttt{║◌│◌◌│H◌║ 7} \\
\texttt{╚═╧══╧══╝} \\
\texttt{```} \\
\texttt{Trenutno igrate krog 6 od največ 20.} \\
\texttt{Skupaj ste zbrali 3 od 12 kazni.} \\
\texttt{Drugi igralec je premaknil predmet na svoji mreži.} \\
\\ 
\\ 
\texttt{Kakšen je vaš naslednji ukaz?} \\
            }
        }
    }
    & \\ \\

    \theutterance \stepcounter{utterance}  
    & \multicolumn{4}{p{0.6\linewidth}}{
        \cellcolor[rgb]{0.9,0.9,0.9}{
            \makecell[{{p{\linewidth}}}]{
                \texttt{\tiny{[P1$\rangle$GM]}}
                \texttt{PREMAKNI SE: T, (7,7)} \\
            }
        }
    }
    & & \\ \\

    \theutterance \stepcounter{utterance}  
    & & & \multicolumn{2}{p{0.3\linewidth}}{
        \cellcolor[rgb]{0.9,0.9,0.9}{
            \makecell[{{p{\linewidth}}}]{
                \texttt{\tiny{[GM$|$GM]}}
                \texttt{PREMAKNI SE: T, (7,7)} \\
            }
        }
    }
    & & \\ \\

    \theutterance \stepcounter{utterance}  
    & & & \multicolumn{2}{p{0.3\linewidth}}{
        \cellcolor[rgb]{0.9,0.9,0.9}{
            \makecell[{{p{\linewidth}}}]{
                \texttt{\tiny{[GM$|$GM]}}
                \texttt{T moved to (7, 7)} \\
            }
        }
    }
    & & \\ \\

    \theutterance \stepcounter{utterance}  
    & & \multicolumn{4}{p{0.6\linewidth}}{
        \cellcolor[rgb]{0.9,0.9,0.9}{
            \makecell[{{p{\linewidth}}}]{
                \texttt{\tiny{[GM$\rangle$P2]}}
                \texttt{Uspešno premaknjen 'I' v (7,4). Vaša posodobljena mreža je videti takole:} \\
\texttt{```} \\
\texttt{ 1234567} \\
\texttt{╔═╤═╤═══╗} \\
\texttt{║H│◌│◌◌◌║ 1} \\
\texttt{║◌├─┤◌◌◌║ 2} \\
\texttt{╟─┤◌├─┬─╢ 3} \\
\texttt{║◌│◌│C│I║ 4} \\
\texttt{╟─┼─┘◌│W║ 5} \\
\texttt{║◌│◌◌┌┴─╢ 6} \\
\texttt{║◌│◌◌│◌T║ 7} \\
\texttt{╚═╧══╧══╝} \\
\texttt{```} \\
\texttt{Trenutno igrate krog 7 od največ 20.} \\
\texttt{Skupaj ste zbrali 3 od 12 kazni.} \\
\texttt{Drugi igralec je premaknil predmet na svoji mreži.} \\
\\ 
\\ 
\texttt{Kakšen je vaš naslednji ukaz?} \\
            }
        }
    }
    & \\ \\

    \theutterance \stepcounter{utterance}  
    & & & \multicolumn{4}{p{0.6\linewidth}}{
        \cellcolor[rgb]{0.9,0.9,0.9}{
            \makecell[{{p{\linewidth}}}]{
                \texttt{\tiny{[GM$\langle$P2]}}
                \texttt{PREMAKNI SE: C, (7, 2)} \\
            }
        }
    }
     \\ \\

    \theutterance \stepcounter{utterance}  
    & & & \multicolumn{2}{p{0.3\linewidth}}{
        \cellcolor[rgb]{0.9,0.9,0.9}{
            \makecell[{{p{\linewidth}}}]{
                \texttt{\tiny{[GM$|$GM]}}
                \texttt{PREMAKNI SE: C, (7, 2)} \\
            }
        }
    }
    & & \\ \\

    \theutterance \stepcounter{utterance}  
    & & & \multicolumn{2}{p{0.3\linewidth}}{
        \cellcolor[rgb]{0.9,0.9,0.9}{
            \makecell[{{p{\linewidth}}}]{
                \texttt{\tiny{[GM$|$GM]}}
                \texttt{C moved to (7, 2)} \\
            }
        }
    }
    & & \\ \\

    \theutterance \stepcounter{utterance}  
    & & \multicolumn{4}{p{0.6\linewidth}}{
        \cellcolor[rgb]{0.9,0.9,0.9}{
            \makecell[{{p{\linewidth}}}]{
                \texttt{\tiny{[P1$\langle$GM]}}
                \texttt{Uspešno premaknjen 'T' v (7,7). Vaša posodobljena mreža je videti takole:} \\
\texttt{```} \\
\texttt{ 1234567} \\
\texttt{╔═╤═╤═══╗} \\
\texttt{║C│◌│◌◌I║ 1} \\
\texttt{║◌├─┤◌◌◌║ 2} \\
\texttt{╟─┤◌├─┬─╢ 3} \\
\texttt{║◌│◌│◌│◌║ 4} \\
\texttt{╟─┼─┘◌│W║ 5} \\
\texttt{║◌│◌◌┌┴─╢ 6} \\
\texttt{║◌│◌◌│HT║ 7} \\
\texttt{╚═╧══╧══╝} \\
\texttt{```} \\
\texttt{Trenutno igrate krog 7 od največ 20.} \\
\texttt{Skupaj ste zbrali 3 od 12 kazni.} \\
\texttt{Drugi igralec je premaknil predmet na svoji mreži.} \\
\\ 
\\ 
\texttt{Kakšen je vaš naslednji ukaz?} \\
            }
        }
    }
    & \\ \\

    \theutterance \stepcounter{utterance}  
    & \multicolumn{4}{p{0.6\linewidth}}{
        \cellcolor[rgb]{0.9,0.9,0.9}{
            \makecell[{{p{\linewidth}}}]{
                \texttt{\tiny{[P1$\rangle$GM]}}
                \texttt{PREMAKNI SE: I, (7,4)} \\
            }
        }
    }
    & & \\ \\

    \theutterance \stepcounter{utterance}  
    & & & \multicolumn{2}{p{0.3\linewidth}}{
        \cellcolor[rgb]{0.9,0.9,0.9}{
            \makecell[{{p{\linewidth}}}]{
                \texttt{\tiny{[GM$|$GM]}}
                \texttt{PREMAKNI SE: I, (7,4)} \\
            }
        }
    }
    & & \\ \\

    \theutterance \stepcounter{utterance}  
    & & & \multicolumn{2}{p{0.3\linewidth}}{
        \cellcolor[rgb]{0.9,0.9,0.9}{
            \makecell[{{p{\linewidth}}}]{
                \texttt{\tiny{[GM$|$GM]}}
                \texttt{I moved to (7, 4)} \\
            }
        }
    }
    & & \\ \\

    \theutterance \stepcounter{utterance}  
    & & \multicolumn{4}{p{0.6\linewidth}}{
        \cellcolor[rgb]{0.9,0.9,0.9}{
            \makecell[{{p{\linewidth}}}]{
                \texttt{\tiny{[GM$\rangle$P2]}}
                \texttt{Uspešno premaknjen 'C' v (7,2). Vaša posodobljena mreža je videti takole:} \\
\texttt{```} \\
\texttt{ 1234567} \\
\texttt{╔═╤═╤═══╗} \\
\texttt{║H│◌│◌◌◌║ 1} \\
\texttt{║◌├─┤◌◌C║ 2} \\
\texttt{╟─┤◌├─┬─╢ 3} \\
\texttt{║◌│◌│◌│I║ 4} \\
\texttt{╟─┼─┘◌│W║ 5} \\
\texttt{║◌│◌◌┌┴─╢ 6} \\
\texttt{║◌│◌◌│◌T║ 7} \\
\texttt{╚═╧══╧══╝} \\
\texttt{```} \\
\texttt{Trenutno igrate krog 8 od največ 20.} \\
\texttt{Skupaj ste zbrali 3 od 12 kazni.} \\
\texttt{Drugi igralec je premaknil predmet na svoji mreži.} \\
\\ 
\\ 
\texttt{Kakšen je vaš naslednji ukaz?} \\
            }
        }
    }
    & \\ \\

    \theutterance \stepcounter{utterance}  
    & & & \multicolumn{4}{p{0.6\linewidth}}{
        \cellcolor[rgb]{0.9,0.9,0.9}{
            \makecell[{{p{\linewidth}}}]{
                \texttt{\tiny{[GM$\langle$P2]}}
                \texttt{PREMAKNI SE: H, (6, 7)} \\
            }
        }
    }
     \\ \\

    \theutterance \stepcounter{utterance}  
    & & & \multicolumn{2}{p{0.3\linewidth}}{
        \cellcolor[rgb]{0.9,0.9,0.9}{
            \makecell[{{p{\linewidth}}}]{
                \texttt{\tiny{[GM$|$GM]}}
                \texttt{PREMAKNI SE: H, (6, 7)} \\
            }
        }
    }
    & & \\ \\

    \theutterance \stepcounter{utterance}  
    & & & \multicolumn{2}{p{0.3\linewidth}}{
        \cellcolor[rgb]{0.9,0.9,0.9}{
            \makecell[{{p{\linewidth}}}]{
                \texttt{\tiny{[GM$|$GM]}}
                \texttt{H moved to (6, 7)} \\
            }
        }
    }
    & & \\ \\

    \theutterance \stepcounter{utterance}  
    & & \multicolumn{4}{p{0.6\linewidth}}{
        \cellcolor[rgb]{0.9,0.9,0.9}{
            \makecell[{{p{\linewidth}}}]{
                \texttt{\tiny{[P1$\langle$GM]}}
                \texttt{Uspešno premaknjen 'I' v (7,4). Vaša posodobljena mreža je videti takole:} \\
\texttt{```} \\
\texttt{ 1234567} \\
\texttt{╔═╤═╤═══╗} \\
\texttt{║C│◌│◌◌◌║ 1} \\
\texttt{║◌├─┤◌◌◌║ 2} \\
\texttt{╟─┤◌├─┬─╢ 3} \\
\texttt{║◌│◌│◌│I║ 4} \\
\texttt{╟─┼─┘◌│W║ 5} \\
\texttt{║◌│◌◌┌┴─╢ 6} \\
\texttt{║◌│◌◌│HT║ 7} \\
\texttt{╚═╧══╧══╝} \\
\texttt{```} \\
\texttt{Trenutno igrate krog 8 od največ 20.} \\
\texttt{Skupaj ste zbrali 3 od 12 kazni.} \\
\texttt{Drugi igralec je premaknil predmet na svoji mreži.} \\
\\ 
\\ 
\texttt{Kakšen je vaš naslednji ukaz?} \\
            }
        }
    }
    & \\ \\

    \theutterance \stepcounter{utterance}  
    & \multicolumn{4}{p{0.6\linewidth}}{
        \cellcolor[rgb]{0.9,0.9,0.9}{
            \makecell[{{p{\linewidth}}}]{
                \texttt{\tiny{[P1$\rangle$GM]}}
                \texttt{PREMAKNI SE: C, (7,2)} \\
            }
        }
    }
    & & \\ \\

    \theutterance \stepcounter{utterance}  
    & & & \multicolumn{2}{p{0.3\linewidth}}{
        \cellcolor[rgb]{0.9,0.9,0.9}{
            \makecell[{{p{\linewidth}}}]{
                \texttt{\tiny{[GM$|$GM]}}
                \texttt{PREMAKNI SE: C, (7,2)} \\
            }
        }
    }
    & & \\ \\

    \theutterance \stepcounter{utterance}  
    & & & \multicolumn{2}{p{0.3\linewidth}}{
        \cellcolor[rgb]{0.9,0.9,0.9}{
            \makecell[{{p{\linewidth}}}]{
                \texttt{\tiny{[GM$|$GM]}}
                \texttt{C moved to (7, 2)} \\
            }
        }
    }
    & & \\ \\

    \theutterance \stepcounter{utterance}  
    & & \multicolumn{4}{p{0.6\linewidth}}{
        \cellcolor[rgb]{0.9,0.9,0.9}{
            \makecell[{{p{\linewidth}}}]{
                \texttt{\tiny{[GM$\rangle$P2]}}
                \texttt{Uspešno premaknjen 'H' v (6,7). Vaša posodobljena mreža je videti takole:} \\
\texttt{```} \\
\texttt{ 1234567} \\
\texttt{╔═╤═╤═══╗} \\
\texttt{║◌│◌│◌◌◌║ 1} \\
\texttt{║◌├─┤◌◌C║ 2} \\
\texttt{╟─┤◌├─┬─╢ 3} \\
\texttt{║◌│◌│◌│I║ 4} \\
\texttt{╟─┼─┘◌│W║ 5} \\
\texttt{║◌│◌◌┌┴─╢ 6} \\
\texttt{║◌│◌◌│HT║ 7} \\
\texttt{╚═╧══╧══╝} \\
\texttt{```} \\
\texttt{Trenutno igrate krog 9 od največ 20.} \\
\texttt{Skupaj ste zbrali 3 od 12 kazni.} \\
\texttt{Drugi igralec je premaknil predmet na svoji mreži.} \\
\\ 
\\ 
\texttt{Kakšen je vaš naslednji ukaz?} \\
            }
        }
    }
    & \\ \\

    \theutterance \stepcounter{utterance}  
    & & & \multicolumn{4}{p{0.6\linewidth}}{
        \cellcolor[rgb]{0.9,0.9,0.9}{
            \makecell[{{p{\linewidth}}}]{
                \texttt{\tiny{[GM$\langle$P2]}}
                \texttt{Vsi moji predmeti so na ciljnih koordinatah: C(7,2), I(7,4), W(7,5), H(6,7), T(7,7). Preverim, če so vsi na pravih mestih:} \\
\texttt{{-} C na (7,2) ✓} \\
\texttt{{-} I na (7,4) ✓} \\
\texttt{{-} W na (7,5) ✓} \\
\texttt{{-} H na (6,7) ✓} \\
\texttt{{-} T na (7,7) ✓} \\
\\ 
\texttt{Vse je poravnano. Zdaj moram preveriti s partnerjem, če je tudi on končal. Poslal bom RECI: končano?RECI: končano?} \\
            }
        }
    }
     \\ \\

    \theutterance \stepcounter{utterance}  
    & & & \multicolumn{2}{p{0.3\linewidth}}{
        \cellcolor[rgb]{0.9,0.9,0.9}{
            \makecell[{{p{\linewidth}}}]{
                \texttt{\tiny{[GM$|$GM]}}
                \texttt{Vsi moji predmeti so na ciljnih koordinatah: C(7,2), I(7,4), W(7,5), H(6,7), T(7,7). Preverim, če so vsi na pravih mestih:} \\
\texttt{{-} C na (7,2) ✓} \\
\texttt{{-} I na (7,4) ✓} \\
\texttt{{-} W na (7,5) ✓} \\
\texttt{{-} H na (6,7) ✓} \\
\texttt{{-} T na (7,7) ✓} \\
\\ 
\texttt{Vse je poravnano. Zdaj moram preveriti s partnerjem, če je tudi on končal. Poslal bom RECI: končano?RECI: končano?} \\
            }
        }
    }
    & & \\ \\

    \theutterance \stepcounter{utterance}  
    & & & \multicolumn{2}{p{0.3\linewidth}}{
        \cellcolor[rgb]{0.9,0.9,0.9}{
            \makecell[{{p{\linewidth}}}]{
                \texttt{\tiny{[GM$|$GM]}}
                \texttt{Invalid format: head is not empty: 'Vsi moji predmeti so na ciljnih koordinatah: C(7,2), I(7,4), W(7,5), H(6,7), T(7,7). Preverim, če so vsi na pravih mestih:} \\
\texttt{{-} C na (7,2) ✓} \\
\texttt{{-} I na (7,4) ✓} \\
\texttt{{-} W na (7,5) ✓} \\
\texttt{{-} H na (6,7) ✓} \\
\texttt{{-} T na (7,7) ✓} \\
\\ 
\texttt{Vse je poravnano. Zdaj moram preveriti s partnerjem, če je tudi on končal. Poslal bom RECI: končano?'} \\
            }
        }
    }
    & & \\ \\

    \theutterance \stepcounter{utterance}  
    & & \multicolumn{4}{p{0.6\linewidth}}{
        \cellcolor[rgb]{0.9,0.9,0.9}{
            \makecell[{{p{\linewidth}}}]{
                \texttt{\tiny{[GM$\rangle$P2]}}
                \texttt{Kazen: Vaše sporočilo ne sme vsebovati ničesar pred ukazom!} \\
\\ 
\texttt{Prepričajte se, da vaš odgovor vsebuje samo `RECI: <SPOROČILO>` ali `PREMAKNI SE: <PREDMET>, (<X>, <Y>)` in nič drugega!} \\
\\ 
\texttt{Skupaj ste zbrali 4 od 12 kazni.} \\
\texttt{Prosim, poskusite znova!} \\
            }
        }
    }
    & \\ \\

    \theutterance \stepcounter{utterance}  
    & & & \multicolumn{4}{p{0.6\linewidth}}{
        \cellcolor[rgb]{0.9,0.9,0.9}{
            \makecell[{{p{\linewidth}}}]{
                \texttt{\tiny{[GM$\langle$P2]}}
                \texttt{RECI: končano?} \\
            }
        }
    }
     \\ \\

    \theutterance \stepcounter{utterance}  
    & & & \multicolumn{2}{p{0.3\linewidth}}{
        \cellcolor[rgb]{0.9,0.9,0.9}{
            \makecell[{{p{\linewidth}}}]{
                \texttt{\tiny{[GM$|$GM]}}
                \texttt{RECI: končano?} \\
            }
        }
    }
    & & \\ \\

    \theutterance \stepcounter{utterance}  
    & & \multicolumn{4}{p{0.6\linewidth}}{
        \cellcolor[rgb]{0.9,0.9,0.9}{
            \makecell[{{p{\linewidth}}}]{
                \texttt{\tiny{[P1$\langle$GM]}}
                \texttt{Uspešno premaknjen 'C' v (7,2). Vaša posodobljena mreža je videti takole:} \\
\texttt{```} \\
\texttt{ 1234567} \\
\texttt{╔═╤═╤═══╗} \\
\texttt{║◌│◌│◌◌◌║ 1} \\
\texttt{║◌├─┤◌◌C║ 2} \\
\texttt{╟─┤◌├─┬─╢ 3} \\
\texttt{║◌│◌│◌│I║ 4} \\
\texttt{╟─┼─┘◌│W║ 5} \\
\texttt{║◌│◌◌┌┴─╢ 6} \\
\texttt{║◌│◌◌│HT║ 7} \\
\texttt{╚═╧══╧══╝} \\
\texttt{```} \\
\texttt{Trenutno igrate krog 9 od največ 20.} \\
\texttt{Skupaj ste zbrali 4 od 12 kazni.} \\
\texttt{Drugi igralec je poslal to sporočilo:} \\
\\ 
\\ 
\texttt{"končano?"} \\
\\ 
\\ 
\texttt{Kakšen je vaš naslednji ukaz?} \\
            }
        }
    }
    & \\ \\

    \theutterance \stepcounter{utterance}  
    & \multicolumn{4}{p{0.6\linewidth}}{
        \cellcolor[rgb]{0.9,0.9,0.9}{
            \makecell[{{p{\linewidth}}}]{
                \texttt{\tiny{[P1$\rangle$GM]}}
                \texttt{RECI: končano!} \\
            }
        }
    }
    & & \\ \\

    \theutterance \stepcounter{utterance}  
    & & & \multicolumn{2}{p{0.3\linewidth}}{
        \cellcolor[rgb]{0.9,0.9,0.9}{
            \makecell[{{p{\linewidth}}}]{
                \texttt{\tiny{[GM$|$GM]}}
                \texttt{RECI: končano!} \\
            }
        }
    }
    & & \\ \\

    \theutterance \stepcounter{utterance}  
    & & & \multicolumn{2}{p{0.3\linewidth}}{
        \cellcolor[rgb]{0.9,0.9,0.9}{
            \makecell[{{p{\linewidth}}}]{
                \texttt{\tiny{[GM$|$GM]}}
                \texttt{true} \\
            }
        }
    }
    & & \\ \\

    \theutterance \stepcounter{utterance}  
    & & & \multicolumn{2}{p{0.3\linewidth}}{
        \cellcolor[rgb]{0.9,0.9,0.9}{
            \makecell[{{p{\linewidth}}}]{
                \texttt{\tiny{[GM$|$GM]}}
                \texttt{Initial states:} \\
\texttt{```} \\
\texttt{Player 1:} \\
\texttt{ 1234567} \\
\texttt{╔═╤═╤═══╗} \\
\texttt{║C│◌│◌TI║ 1} \\
\texttt{║◌├─┤W◌◌║ 2} \\
\texttt{╟─┤◌├─┬─╢ 3} \\
\texttt{║◌│◌│◌│◌║ 4} \\
\texttt{╟─┼─┘◌│◌║ 5} \\
\texttt{║◌│◌H┌┴─╢ 6} \\
\texttt{║◌│◌◌│◌◌║ 7} \\
\texttt{╚═╧══╧══╝} \\
\\ 
\\ 
\texttt{Player 2:} \\
\texttt{ 1234567} \\
\texttt{╔═╤═╤═══╗} \\
\texttt{║H│◌│◌◌◌║ 1} \\
\texttt{║◌├─┤◌◌◌║ 2} \\
\texttt{╟─┤◌├─┬─╢ 3} \\
\texttt{║◌│◌│C│◌║ 4} \\
\texttt{╟─┼─┘◌│◌║ 5} \\
\texttt{║◌│◌◌┌┴─╢ 6} \\
\texttt{║T│WI│◌◌║ 7} \\
\texttt{╚═╧══╧══╝} \\
\\ 
\texttt{```} \\
            }
        }
    }
    & & \\ \\

    \theutterance \stepcounter{utterance}  
    & & & \multicolumn{2}{p{0.3\linewidth}}{
        \cellcolor[rgb]{0.9,0.9,0.9}{
            \makecell[{{p{\linewidth}}}]{
                \texttt{\tiny{[GM$|$GM]}}
                \texttt{End states:} \\
\texttt{```} \\
\texttt{Player 1:} \\
\texttt{ 1234567} \\
\texttt{╔═╤═╤═══╗} \\
\texttt{║◌│◌│◌◌◌║ 1} \\
\texttt{║◌├─┤◌◌C║ 2} \\
\texttt{╟─┤◌├─┬─╢ 3} \\
\texttt{║◌│◌│◌│I║ 4} \\
\texttt{╟─┼─┘◌│W║ 5} \\
\texttt{║◌│◌◌┌┴─╢ 6} \\
\texttt{║◌│◌◌│HT║ 7} \\
\texttt{╚═╧══╧══╝} \\
\\ 
\\ 
\texttt{Player 2:} \\
\texttt{ 1234567} \\
\texttt{╔═╤═╤═══╗} \\
\texttt{║◌│◌│◌◌◌║ 1} \\
\texttt{║◌├─┤◌◌C║ 2} \\
\texttt{╟─┤◌├─┬─╢ 3} \\
\texttt{║◌│◌│◌│I║ 4} \\
\texttt{╟─┼─┘◌│W║ 5} \\
\texttt{║◌│◌◌┌┴─╢ 6} \\
\texttt{║◌│◌◌│HT║ 7} \\
\texttt{╚═╧══╧══╝} \\
\\ 
\texttt{```} \\
            }
        }
    }
    & & \\ \\

    \theutterance \stepcounter{utterance}  
    & & & \multicolumn{2}{p{0.3\linewidth}}{
        \cellcolor[rgb]{0.9,0.9,0.9}{
            \makecell[{{p{\linewidth}}}]{
                \texttt{\tiny{[GM$|$GM]}}
                \texttt{* success: True} \\
\texttt{* lose: False} \\
\texttt{* aborted: False} \\
\texttt{{-}{-}{-}{-}{-}{-}{-}} \\
\texttt{* Shifts: 8.00} \\
\texttt{* Max Shifts: 8.00} \\
\texttt{* Min Shifts: 4.00} \\
\texttt{* End Distance Sum: 0.00} \\
\texttt{* Init Distance Sum: 30.73} \\
\texttt{* Expected Distance Sum: 20.95} \\
\texttt{* Penalties: 4.00} \\
\texttt{* Max Penalties: 12.00} \\
\texttt{* Rounds: 9.00} \\
\texttt{* Max Rounds: 20.00} \\
\texttt{* Object Count: 5.00} \\
            }
        }
    }
    & & \\ \\

\end{supertabular}
}

\end{document}
